\scanpage{99}
% archive foliation: 98
trouve que la catégorie des cartes \inserted{\uline{orientées}} $\mathcal C_0$-épinglées est rigide -- donc équivalente à une catégorie discrète. \deleted{\unsure{Donc}} Dans le cas où $\mathcal C_0$ est fini, cette catégorie est \unsure{elle} $\simeq$ finie. Elle est vide sauf si
\[
\sum_s p(s)=\sum_a q(a)=\sum_f r(f),
\]
et si cette condition est vérifiée, soit \deleted{$\unsure{d}$}\inserted{$\delta$} la \deleted{\unsure{valeur}}\inserted{valeur} commune des trois -- alors le \deleted{\unsure{nombre des}}\inserted{cardinal de l'} classes d'isomorphie des $\mathcal C_0$-cartes est égal à \unsure{l'ens.} des classes d'iso. de $\mathcal G_2^\circ$-ens $E$ ``admissibles'' de cardinal $\delta$, ayant les \unsure{mêmes} \deleted{\unsure{invariants}} nombres de multiplicité que $\mathcal C_0$, \uline{multiplié} par le cardinal du groupe des autom. de $\mathcal C_0$, qui est égal à
\[
\prod_{n\in\mathbb N^*}(a_n!\,b_n!\,c_n!)
\]
où $a_n(b_n,c_n)$ est le nb de sommets (arêtes, faces) de degré $n$.

\margin{Définir ici $C_{\mathcal C_0}$ et $C_{\mathcal C_0}^{\alpha}$, et l'opération de $\gamma=\gamma_{\mathcal C_0}=\operatorname{Aut}(\mathcal C_0)$ dessus. -- \uline{NB} $C_{\mathcal C_0}^{\alpha}/\gamma\;\unsure{-|\longrightarrow}\;$ l'ens des classes d'iso. de cartes de type $(\mathcal C_0,\alpha)$\ldots}
